\chapter{Predicted Provision Mapping}\label{app:mapping}

The table below is the mapping this project produces: every pair of provisions
that the best-scoring method flags as related at its calibrated
threshold, with the similarity score that produced the flag. It is included
because it is the artefact a compliance practitioner would actually use, and
because a comparative mapping study is expected to show its mapping rather than
only summary statistics of it.

\medskip
\noindent\textbf{This is a predicted mapping, not a verified one.} Every row is
the output of an automated method whose precision on the annotated reference set
was well below one; the measured value and its confidence interval are printed
in Table~\ref{tab:calibrated} and in the caption of Figure~\ref{fig:coverage}. A
substantial fraction of the rows below are therefore wrong. The table is a
shortlist: the pairs an analyst would examine first, in the order the method
ranks them. The measured statements in this thesis are those in
Chapter~\ref{ch:results}, which carry intervals and are scored against annotated
labels; nothing in this appendix is.

\medskip
\noindent Rows are grouped by \standardA{} provision and ordered within a group
by descending score. The table is generated directly from the evaluation
outputs by the \texttt{coverage} module, so it cannot drift from the data
it summarises.

\bigskip
% Generated by src.evaluation.coverage -- do not edit by hand.
\begin{longtable}{@{}lll@{}}
\toprule
EN 303 645 & EN 304 223 & Score \\
\midrule
\endfirsthead
\toprule
EN 303 645 & EN 304 223 & Score \\
\midrule
\endhead
\bottomrule
\endfoot
5.0-1 & 5.1.3-1 & 0.751 \\
5.0-1 & 5.2.3-2 & 0.749 \\
5.0-1 & 5.2.1-1 & 0.742 \\
5.0-1 & 5.1.2-1 & 0.741 \\
5.0-1 & 5.1.3-1.1 & 0.734 \\
5.0-1 & 5.1.2-7 & 0.727 \\
5.0-1 & 5.2.3-2.1 & 0.725 \\
5.0-1 & 5.1.3-2 & 0.725 \\
5.0-1 & 5.1.2-3 & 0.723 \\
5.0-1 & 5.2.4-1.1 & 0.722 \\
5.0-1 & 5.2.2-4 & 0.722 \\
5.1-1 & 5.4.1-1 & 0.733 \\
5.1-1 & 5.2.1-2 & 0.732 \\
5.1-1 & 5.2.1-4 & 0.729 \\
5.1-1 & 5.2.5-1 & 0.727 \\
5.1-1 & 5.2.3-1 & 0.720 \\
5.1-1 & 5.1.3-3 & 0.720 \\
5.1-1 & 5.1.2-7 & 0.719 \\
5.1-2 & 5.2.1-4 & 0.758 \\
5.1-2 & 5.1.3-1 & 0.757 \\
5.1-2 & 5.1.3-3 & 0.754 \\
5.1-2 & 5.2.1-4.1 & 0.753 \\
5.1-2 & 5.2.3-1 & 0.752 \\
5.1-2 & 5.2.1-2 & 0.751 \\
5.1-2 & 5.2.3-2.1 & 0.751 \\
5.1-2 & 5.1.2-4 & 0.749 \\
5.1-2 & 5.1.2-2 & 0.749 \\
5.1-2 & 5.4.1-1 & 0.745 \\
5.1-2 & 5.2.2-4 & 0.742 \\
5.1-2 & 5.2.5-1 & 0.742 \\
5.1-2 & 5.2.5-2.1 & 0.740 \\
5.1-2 & 5.4.1-1.1 & 0.739 \\
5.1-2 & 5.1.2-7 & 0.738 \\
5.1-2 & 5.1.4-3 & 0.737 \\
5.1-2 & 5.2.1-1 & 0.737 \\
5.1-2 & 5.2.2-1 & 0.735 \\
5.1-2 & 5.1.2-6 & 0.734 \\
5.1-2 & 5.1.2-1 & 0.732 \\
5.1-2 & 5.3.1-3 & 0.732 \\
5.1-2 & 5.1.3-1.1 & 0.731 \\
5.1-2 & 5.1.3-2 & 0.730 \\
5.1-2 & 5.1.1-2.2 & 0.730 \\
5.1-2 & 5.2.2-2 & 0.729 \\
5.1-2 & 5.3.1-2.2 & 0.728 \\
5.1-2 & 5.2.3-2 & 0.728 \\
5.1-2 & 5.2.1-4.2 & 0.727 \\
5.1-2 & 5.1.1-2 & 0.725 \\
5.1-2 & 5.1.1-1 & 0.725 \\
5.1-3 & 5.2.1-2 & 0.760 \\
5.1-3 & 5.2.2-1 & 0.746 \\
5.1-3 & 5.2.3-1 & 0.746 \\
5.1-3 & 5.2.5-1 & 0.744 \\
5.1-3 & 5.1.3-3 & 0.744 \\
5.1-3 & 5.2.1-4 & 0.742 \\
5.1-3 & 5.2.1-4.1 & 0.737 \\
5.1-3 & 5.2.2-4 & 0.736 \\
5.1-3 & 5.1.2-7 & 0.736 \\
5.1-3 & 5.1.3-1 & 0.732 \\
5.1-3 & 5.1.2-2 & 0.730 \\
5.1-3 & 5.2.3-2.1 & 0.730 \\
5.1-3 & 5.4.1-1 & 0.729 \\
5.1-3 & 5.2.5-2.1 & 0.728 \\
5.1-3 & 5.1.2-1 & 0.727 \\
5.1-3 & 5.1.2-4 & 0.724 \\
5.1-3 & 5.3.1-3 & 0.724 \\
5.1-3 & 5.2.3-2 & 0.723 \\
5.1-3 & 5.2.4-1.2 & 0.722 \\
5.1-4 & 5.4.1-1 & 0.746 \\
5.1-4 & 5.3.1-2.2 & 0.743 \\
5.1-4 & 5.2.2-4 & 0.740 \\
5.1-4 & 5.2.1-4 & 0.728 \\
5.1-4 & 5.2.5-1 & 0.724 \\
5.1-4 & 5.2.1-2 & 0.721 \\
5.1-5 & 5.1.2-2 & 0.745 \\
5.1-5 & 5.2.1-4 & 0.744 \\
5.1-5 & 5.2.1-2 & 0.743 \\
5.1-5 & 5.2.2-1 & 0.742 \\
5.1-5 & 5.2.2-2 & 0.741 \\
5.1-5 & 5.1.3-1 & 0.740 \\
5.1-5 & 5.2.5-1 & 0.736 \\
5.1-5 & 5.1.3-3 & 0.726 \\
5.1-5 & 5.2.2-4 & 0.726 \\
5.1-5 & 5.2.3-1 & 0.723 \\
5.1-5 & 5.1.2-1 & 0.723 \\
5.10-1 & 5.4.2-3 & 0.776 \\
5.10-1 & 5.4.2-4 & 0.746 \\
5.10-1 & 5.4.2-2 & 0.745 \\
5.11-1 & 5.5.1-2 & 0.733 \\
5.11-1 & 5.3.1-1 & 0.730 \\
5.11-2 & 5.5.1-2 & 0.761 \\
5.11-2 & 5.3.1-1 & 0.721 \\
5.11-2 & 5.5.1-1 & 0.720 \\
5.11-3 & 5.5.1-2 & 0.770 \\
5.11-3 & 5.3.1-1 & 0.768 \\
5.11-3 & 5.3.1-2.2 & 0.732 \\
5.11-3 & 5.4.1-1 & 0.729 \\
5.11-3 & 5.5.1-1 & 0.727 \\
5.11-3 & 5.3.1-2 & 0.723 \\
5.11-3 & 5.1.4-5 & 0.719 \\
5.11-4 & 5.5.1-2 & 0.773 \\
5.11-4 & 5.3.1-1 & 0.742 \\
5.11-4 & 5.5.1-1 & 0.726 \\
5.11-4 & 5.2.2-6 & 0.721 \\
5.12-1 & 5.3.1-2.2 & 0.751 \\
5.12-1 & 5.4.1-1 & 0.732 \\
5.12-1 & 5.1.4-2 & 0.725 \\
5.12-1 & 5.2.2-4 & 0.720 \\
5.12-2 & 5.3.1-2.2 & 0.765 \\
5.12-2 & 5.4.1-1 & 0.757 \\
5.12-2 & 5.2.1-1 & 0.754 \\
5.12-2 & 5.3.1-3 & 0.747 \\
5.12-2 & 5.1.3-2 & 0.736 \\
5.12-2 & 5.3.1-2 & 0.733 \\
5.12-2 & 5.2.2-4 & 0.729 \\
5.12-2 & 5.4.1-2 & 0.727 \\
5.12-2 & 5.1.3-1.1 & 0.725 \\
5.12-2 & 5.1.3-1 & 0.723 \\
5.12-3 & 5.3.1-2.2 & 0.800 \\
5.12-3 & 5.4.1-1 & 0.778 \\
5.12-3 & 5.1.3-2 & 0.769 \\
5.12-3 & 5.3.1-2 & 0.765 \\
5.12-3 & 5.3.1-3 & 0.763 \\
5.12-3 & 5.4.2-4 & 0.758 \\
5.12-3 & 5.4.2-3 & 0.754 \\
5.12-3 & 5.2.5-1 & 0.752 \\
5.12-3 & 5.3.1-1 & 0.752 \\
5.12-3 & 5.1.3-3 & 0.750 \\
5.12-3 & 5.4.1-2 & 0.749 \\
5.12-3 & 5.2.2-4 & 0.748 \\
5.12-3 & 5.1.2-7 & 0.747 \\
5.12-3 & 5.1.3-1 & 0.742 \\
5.12-3 & 5.2.3-2.2 & 0.739 \\
5.12-3 & 5.2.1-4.1 & 0.737 \\
5.12-3 & 5.4.1-1.1 & 0.736 \\
5.12-3 & 5.1.4-2 & 0.735 \\
5.12-3 & 5.2.1-1 & 0.734 \\
5.12-3 & 5.2.4-1.1 & 0.731 \\
5.12-3 & 5.1.3-4 & 0.730 \\
5.12-3 & 5.1.3-1.1 & 0.729 \\
5.12-3 & 5.2.3-1 & 0.729 \\
5.12-3 & 5.1.1-2 & 0.729 \\
5.12-3 & 5.3.1-2.1 & 0.729 \\
5.12-3 & 5.2.5-2.1 & 0.728 \\
5.12-3 & 5.2.5-4.1 & 0.728 \\
5.12-3 & 5.4.2-2 & 0.728 \\
5.12-3 & 5.1.4-1 & 0.728 \\
5.12-3 & 5.1.2-1 & 0.727 \\
5.12-3 & 5.1.1-1 & 0.724 \\
5.12-3 & 5.1.2-4 & 0.723 \\
5.12-3 & 5.2.1-2 & 0.721 \\
5.12-3 & 5.2.3-2 & 0.721 \\
5.13-1A & 5.2.1-4.1 & 0.751 \\
5.13-1A & 5.2.1-4 & 0.736 \\
5.13-1A & 5.2.5-1 & 0.721 \\
5.13-1B & 5.2.1-4.1 & 0.766 \\
5.13-1B & 5.2.5-1 & 0.746 \\
5.13-1B & 5.2.1-4 & 0.738 \\
5.13-1B & 5.1.2-2 & 0.730 \\
5.13-1B & 5.1.3-1 & 0.726 \\
5.13-1B & 5.2.1-1 & 0.726 \\
5.13-1B & 5.1.4-4 & 0.726 \\
5.2-1 & 5.2.2-4 & 0.871 \\
5.2-1 & 5.4.1-1 & 0.778 \\
5.2-1 & 5.1.3-1 & 0.767 \\
5.2-1 & 5.2.5-1 & 0.767 \\
5.2-1 & 5.4.1-1.1 & 0.766 \\
5.2-1 & 5.2.1-1 & 0.763 \\
5.2-1 & 5.2.3-1 & 0.763 \\
5.2-1 & 5.1.3-2 & 0.759 \\
5.2-1 & 5.3.1-3 & 0.759 \\
5.2-1 & 5.1.3-3 & 0.758 \\
5.2-1 & 5.3.1-2.2 & 0.757 \\
5.2-1 & 5.1.2-4 & 0.748 \\
5.2-1 & 5.1.2-7 & 0.747 \\
5.2-1 & 5.1.2-1 & 0.747 \\
5.2-1 & 5.1.3-1.1 & 0.746 \\
5.2-1 & 5.4.1-2 & 0.741 \\
5.2-1 & 5.1.1-2 & 0.740 \\
5.2-1 & 5.2.1-2 & 0.740 \\
5.2-1 & 5.2.3-2 & 0.738 \\
5.2-1 & 5.2.5-2.1 & 0.737 \\
5.2-1 & 5.1.1-1 & 0.736 \\
5.2-1 & 5.2.3-2.1 & 0.734 \\
5.2-1 & 5.1.2-2 & 0.734 \\
5.2-1 & 5.2.4-2 & 0.731 \\
5.2-1 & 5.2.4-1.1 & 0.730 \\
5.2-1 & 5.3.1-1 & 0.730 \\
5.2-1 & 5.2.2-5 & 0.729 \\
5.2-1 & 5.1.1-2.2 & 0.725 \\
5.2-1 & 5.1.2-3 & 0.725 \\
5.2-1 & 5.2.1-4.1 & 0.724 \\
5.2-1 & 5.2.1-4 & 0.723 \\
5.2-1 & 5.1.2-5.1 & 0.722 \\
5.2-1 & 5.2.4-1 & 0.720 \\
5.2-2 & 5.2.2-4 & 0.781 \\
5.2-2 & 5.4.1-1 & 0.728 \\
5.2-3 & 5.4.1-1 & 0.787 \\
5.2-3 & 5.2.2-4 & 0.786 \\
5.2-3 & 5.1.3-1 & 0.774 \\
5.2-3 & 5.3.1-3 & 0.767 \\
5.2-3 & 5.4.2-4 & 0.766 \\
5.2-3 & 5.4.2-3 & 0.764 \\
5.2-3 & 5.1.1-2 & 0.760 \\
5.2-3 & 5.4.1-2 & 0.759 \\
5.2-3 & 5.4.1-1.1 & 0.756 \\
5.2-3 & 5.3.1-2.2 & 0.755 \\
5.2-3 & 5.1.3-4 & 0.754 \\
5.2-3 & 5.2.5-1 & 0.751 \\
5.2-3 & 5.2.1-1 & 0.744 \\
5.2-3 & 5.1.3-2 & 0.742 \\
5.2-3 & 5.2.3-1 & 0.742 \\
5.2-3 & 5.1.3-3 & 0.741 \\
5.2-3 & 5.1.1-1 & 0.738 \\
5.2-3 & 5.1.3-1.1 & 0.738 \\
5.2-3 & 5.2.1-4.1 & 0.734 \\
5.2-3 & 5.1.1-2.2 & 0.734 \\
5.2-3 & 5.1.2-7 & 0.732 \\
5.2-3 & 5.1.2-2 & 0.732 \\
5.2-3 & 5.2.1-2 & 0.730 \\
5.2-3 & 5.1.2-4 & 0.728 \\
5.2-3 & 5.2.5-2.1 & 0.724 \\
5.2-3 & 5.4.2-2 & 0.723 \\
5.3-1 & 5.4.1-1 & 0.790 \\
5.3-1 & 5.3.1-2.2 & 0.776 \\
5.3-1 & 5.4.1-2 & 0.760 \\
5.3-1 & 5.2.3-1 & 0.758 \\
5.3-1 & 5.4.1-1.1 & 0.753 \\
5.3-1 & 5.2.2-4 & 0.748 \\
5.3-1 & 5.2.5-1 & 0.722 \\
5.3-1 & 5.1.3-1 & 0.721 \\
5.3-1 & 5.2.1-1 & 0.720 \\
5.3-1 & 5.1.2-4 & 0.719 \\
5.3-10 & 5.4.1-1 & 0.786 \\
5.3-10 & 5.2.3-1 & 0.774 \\
5.3-10 & 5.3.1-2.2 & 0.762 \\
5.3-10 & 5.1.3-3 & 0.758 \\
5.3-10 & 5.2.5-1 & 0.755 \\
5.3-10 & 5.1.2-4 & 0.752 \\
5.3-10 & 5.4.1-1.1 & 0.751 \\
5.3-10 & 5.4.1-2 & 0.749 \\
5.3-10 & 5.2.4-1.2 & 0.745 \\
5.3-10 & 5.2.1-4.1 & 0.743 \\
5.3-10 & 5.2.2-4 & 0.743 \\
5.3-10 & 5.1.3-1 & 0.743 \\
5.3-10 & 5.2.1-2 & 0.738 \\
5.3-10 & 5.1.2-2 & 0.737 \\
5.3-10 & 5.1.2-7 & 0.736 \\
5.3-10 & 5.2.3-2.1 & 0.735 \\
5.3-10 & 5.1.1-1 & 0.734 \\
5.3-10 & 5.2.1-1 & 0.731 \\
5.3-10 & 5.1.3-2 & 0.728 \\
5.3-10 & 5.1.1-2 & 0.726 \\
5.3-10 & 5.2.1-4 & 0.726 \\
5.3-10 & 5.1.4-3 & 0.725 \\
5.3-10 & 5.1.3-1.1 & 0.724 \\
5.3-10 & 5.2.3-2 & 0.722 \\
5.3-10 & 5.3.1-3 & 0.722 \\
5.3-10 & 5.2.1-4.2 & 0.721 \\
5.3-10 & 5.5.1-1 & 0.721 \\
5.3-11 & 5.3.1-2.2 & 0.814 \\
5.3-11 & 5.4.1-1 & 0.799 \\
5.3-11 & 5.4.1-2 & 0.749 \\
5.3-11 & 5.2.2-4 & 0.747 \\
5.3-11 & 5.1.3-2 & 0.740 \\
5.3-11 & 5.1.1-2.1 & 0.737 \\
5.3-11 & 5.4.1-1.1 & 0.735 \\
5.3-12 & 5.3.1-2.2 & 0.759 \\
5.3-12 & 5.2.3-4 & 0.754 \\
5.3-12 & 5.4.1-1 & 0.737 \\
5.3-13 & 5.3.1-2.2 & 0.814 \\
5.3-13 & 5.4.1-1 & 0.802 \\
5.3-13 & 5.2.2-4 & 0.788 \\
5.3-13 & 5.3.1-1 & 0.772 \\
5.3-13 & 5.2.3-4 & 0.767 \\
5.3-13 & 5.2.3-2.2 & 0.766 \\
5.3-13 & 5.3.1-2 & 0.764 \\
5.3-13 & 5.3.1-3 & 0.751 \\
5.3-13 & 5.4.1-1.1 & 0.750 \\
5.3-13 & 5.2.3-1 & 0.745 \\
5.3-13 & 5.2.5-1 & 0.735 \\
5.3-13 & 5.1.3-2 & 0.734 \\
5.3-13 & 5.2.3-2 & 0.733 \\
5.3-13 & 5.1.2-7 & 0.732 \\
5.3-13 & 5.4.1-2 & 0.729 \\
5.3-13 & 5.1.3-1 & 0.729 \\
5.3-13 & 5.2.1-1 & 0.727 \\
5.3-13 & 5.1.2-4 & 0.726 \\
5.3-13 & 5.1.2-2 & 0.724 \\
5.3-13 & 5.4.1-3 & 0.722 \\
5.3-13 & 5.1.4-5 & 0.721 \\
5.3-13 & 5.2.4-1 & 0.720 \\
5.3-13 & 5.1.2-3 & 0.720 \\
5.3-14 & 5.3.1-2.2 & 0.813 \\
5.3-14 & 5.2.2-4 & 0.790 \\
5.3-14 & 5.4.1-1 & 0.790 \\
5.3-14 & 5.3.1-1 & 0.771 \\
5.3-14 & 5.4.1-1.1 & 0.766 \\
5.3-14 & 5.3.1-3 & 0.766 \\
5.3-14 & 5.3.1-2 & 0.765 \\
5.3-14 & 5.2.3-2.2 & 0.760 \\
5.3-14 & 5.2.3-4 & 0.754 \\
5.3-14 & 5.1.3-2 & 0.745 \\
5.3-14 & 5.2.3-1 & 0.743 \\
5.3-14 & 5.2.3-2 & 0.740 \\
5.3-14 & 5.2.1-1 & 0.736 \\
5.3-14 & 5.4.1-2 & 0.733 \\
5.3-14 & 5.1.2-4 & 0.730 \\
5.3-14 & 5.1.2-7 & 0.730 \\
5.3-14 & 5.2.2-6 & 0.727 \\
5.3-14 & 5.1.4-5 & 0.726 \\
5.3-14 & 5.2.4-1 & 0.725 \\
5.3-14 & 5.2.3-2.1 & 0.724 \\
5.3-14 & 5.1.2-2 & 0.723 \\
5.3-14 & 5.1.3-3 & 0.723 \\
5.3-14 & 5.1.4-2 & 0.723 \\
5.3-14 & 5.4.1-3 & 0.720 \\
5.3-14 & 5.2.5-1 & 0.720 \\
5.3-14 & 5.1.3-1 & 0.720 \\
5.3-15B & 5.4.1-1.1 & 0.749 \\
5.3-15B & 5.4.1-1 & 0.749 \\
5.3-15B & 5.3.1-2.2 & 0.732 \\
5.3-15B & 5.4.1-2 & 0.732 \\
5.3-15B & 5.2.2-4 & 0.723 \\
5.3-16 & 5.2.5-1 & 0.733 \\
5.3-16 & 5.2.3-4 & 0.731 \\
5.3-16 & 5.3.1-1 & 0.730 \\
5.3-16 & 5.4.1-3 & 0.725 \\
5.3-16 & 5.4.1-1 & 0.723 \\
5.3-16 & 5.2.3-2.1 & 0.721 \\
5.3-16 & 5.3.1-2.2 & 0.719 \\
5.3-16 & 5.2.1-1 & 0.719 \\
5.3-2 & 5.4.1-1 & 0.781 \\
5.3-2 & 5.4.1-1.1 & 0.770 \\
5.3-2 & 5.3.1-2.2 & 0.765 \\
5.3-2 & 5.2.3-1 & 0.762 \\
5.3-2 & 5.2.2-4 & 0.759 \\
5.3-2 & 5.2.5-1 & 0.740 \\
5.3-2 & 5.4.1-2 & 0.739 \\
5.3-2 & 5.2.3-2.1 & 0.738 \\
5.3-2 & 5.2.3-2 & 0.734 \\
5.3-2 & 5.1.2-4 & 0.732 \\
5.3-2 & 5.1.3-1.1 & 0.730 \\
5.3-2 & 5.2.1-2 & 0.730 \\
5.3-2 & 5.1.2-1 & 0.729 \\
5.3-2 & 5.1.3-1 & 0.729 \\
5.3-2 & 5.1.1-2 & 0.729 \\
5.3-2 & 5.2.1-1 & 0.727 \\
5.3-2 & 5.1.1-1 & 0.726 \\
5.3-2 & 5.1.3-3 & 0.724 \\
5.3-2 & 5.1.3-2 & 0.720 \\
5.3-3 & 5.3.1-2.2 & 0.775 \\
5.3-3 & 5.4.1-1 & 0.774 \\
5.3-3 & 5.2.2-4 & 0.729 \\
5.3-4A & 5.4.1-1 & 0.750 \\
5.3-4A & 5.2.1-3.1 & 0.739 \\
5.3-4A & 5.3.1-2.2 & 0.729 \\
5.3-4A & 5.2.2-4 & 0.725 \\
5.3-4A & 5.4.1-2 & 0.722 \\
5.3-4B & 5.4.1-1 & 0.788 \\
5.3-4B & 5.3.1-2.2 & 0.774 \\
5.3-4B & 5.4.1-1.1 & 0.768 \\
5.3-4B & 5.4.1-2 & 0.765 \\
5.3-4B & 5.3.1-3 & 0.753 \\
5.3-4B & 5.2.2-4 & 0.747 \\
5.3-4B & 5.2.1-3.1 & 0.740 \\
5.3-4B & 5.1.3-1 & 0.734 \\
5.3-4B & 5.1.3-1.1 & 0.731 \\
5.3-4B & 5.2.3-1 & 0.730 \\
5.3-4B & 5.1.2-2 & 0.727 \\
5.3-4B & 5.2.2-5 & 0.726 \\
5.3-4B & 5.1.3-2 & 0.726 \\
5.3-4B & 5.1.2-4 & 0.724 \\
5.3-4B & 5.2.1-4.1 & 0.722 \\
5.3-4B & 5.1.1-1 & 0.720 \\
5.3-5 & 5.4.1-1 & 0.761 \\
5.3-5 & 5.4.1-2 & 0.754 \\
5.3-5 & 5.3.1-2.2 & 0.748 \\
5.3-5 & 5.2.2-4 & 0.724 \\
5.3-5 & 5.4.1-1.1 & 0.724 \\
5.3-6A & 5.4.1-1 & 0.768 \\
5.3-6A & 5.3.1-2.2 & 0.762 \\
5.3-6A & 5.4.1-2 & 0.729 \\
5.3-6A & 5.4.1-1.1 & 0.721 \\
5.3-6B & 5.3.1-2.2 & 0.769 \\
5.3-6B & 5.4.1-1 & 0.760 \\
5.3-6B & 5.3.1-1 & 0.721 \\
5.3-6B & 5.2.3-4 & 0.719 \\
5.3-7 & 5.4.1-1 & 0.777 \\
5.3-7 & 5.2.3-1 & 0.760 \\
5.3-7 & 5.3.1-2.2 & 0.758 \\
5.3-7 & 5.4.1-1.1 & 0.740 \\
5.3-7 & 5.2.5-1 & 0.738 \\
5.3-7 & 5.2.2-4 & 0.737 \\
5.3-7 & 5.4.1-2 & 0.736 \\
5.3-7 & 5.1.3-1 & 0.731 \\
5.3-7 & 5.2.4-1.2 & 0.727 \\
5.3-7 & 5.2.1-2 & 0.724 \\
5.3-7 & 5.1.3-1.1 & 0.721 \\
5.3-8 & 5.4.1-1 & 0.802 \\
5.3-8 & 5.3.1-2.2 & 0.790 \\
5.3-8 & 5.2.2-4 & 0.754 \\
5.3-8 & 5.4.1-1.1 & 0.744 \\
5.3-8 & 5.4.1-2 & 0.744 \\
5.3-8 & 5.1.3-2 & 0.738 \\
5.3-8 & 5.1.1-2.1 & 0.729 \\
5.3-8 & 5.3.1-3 & 0.720 \\
5.3-9 & 5.4.1-2 & 0.736 \\
5.3-9 & 5.4.1-1 & 0.731 \\
5.3-9 & 5.2.4-1.2 & 0.731 \\
5.4-1 & 5.2.1-4 & 0.772 \\
5.4-1 & 5.2.1-2 & 0.746 \\
5.4-1 & 5.2.3-1 & 0.736 \\
5.4-1 & 5.2.1-4.2 & 0.733 \\
5.4-1 & 5.2.1-1 & 0.732 \\
5.4-1 & 5.2.2-4 & 0.729 \\
5.4-1 & 5.4.1-1 & 0.724 \\
5.4-1 & 5.1.3-1 & 0.724 \\
5.4-1 & 5.2.5-1 & 0.723 \\
5.4-2 & 5.1.2-2 & 0.758 \\
5.4-2 & 5.2.1-2 & 0.758 \\
5.4-2 & 5.2.3-1 & 0.756 \\
5.4-2 & 5.2.1-4 & 0.750 \\
5.4-2 & 5.1.3-1 & 0.745 \\
5.4-2 & 5.2.2-4 & 0.740 \\
5.4-2 & 5.1.1-2.2 & 0.740 \\
5.4-2 & 5.4.1-1 & 0.740 \\
5.4-2 & 5.2.1-4.1 & 0.739 \\
5.4-2 & 5.2.1-4.2 & 0.737 \\
5.4-2 & 5.2.5-1 & 0.737 \\
5.4-2 & 5.2.1-1 & 0.734 \\
5.4-2 & 5.2.2-1 & 0.732 \\
5.4-2 & 5.1.3-3 & 0.730 \\
5.4-2 & 5.4.1-1.1 & 0.726 \\
5.4-2 & 5.1.1-2 & 0.724 \\
5.4-2 & 5.1.2-1 & 0.723 \\
5.4-2 & 5.2.2-3 & 0.723 \\
5.4-2 & 5.4.2-3 & 0.722 \\
5.4-2 & 5.2.5-2.1 & 0.722 \\
5.4-2 & 5.3.1-3 & 0.721 \\
5.4-2 & 5.2.3-2.1 & 0.721 \\
5.4-2 & 5.1.2-4 & 0.720 \\
5.4-2 & 5.5.1-1 & 0.720 \\
5.4-3 & 5.2.3-1 & 0.751 \\
5.4-3 & 5.2.5-1 & 0.751 \\
5.4-3 & 5.2.1-2 & 0.748 \\
5.4-3 & 5.1.1-2.2 & 0.739 \\
5.4-3 & 5.2.2-4 & 0.737 \\
5.4-3 & 5.1.3-1 & 0.728 \\
5.4-3 & 5.4.1-1 & 0.723 \\
5.4-3 & 5.1.2-7 & 0.722 \\
5.4-3 & 5.2.1-4 & 0.720 \\
5.4-4 & 5.2.3-1 & 0.790 \\
5.4-4 & 5.4.1-1 & 0.783 \\
5.4-4 & 5.1.3-1 & 0.773 \\
5.4-4 & 5.2.1-2 & 0.771 \\
5.4-4 & 5.4.1-1.1 & 0.764 \\
5.4-4 & 5.3.1-2.2 & 0.764 \\
5.4-4 & 5.1.2-2 & 0.762 \\
5.4-4 & 5.2.5-1 & 0.762 \\
5.4-4 & 5.1.2-4 & 0.757 \\
5.4-4 & 5.2.1-4.1 & 0.757 \\
5.4-4 & 5.2.2-4 & 0.756 \\
5.4-4 & 5.2.1-4 & 0.753 \\
5.4-4 & 5.2.3-2.1 & 0.751 \\
5.4-4 & 5.1.3-1.1 & 0.750 \\
5.4-4 & 5.1.1-2.2 & 0.748 \\
5.4-4 & 5.1.1-1 & 0.748 \\
5.4-4 & 5.1.3-3 & 0.748 \\
5.4-4 & 5.4.1-2 & 0.748 \\
5.4-4 & 5.2.1-1 & 0.746 \\
5.4-4 & 5.2.5-2.1 & 0.744 \\
5.4-4 & 5.2.3-2 & 0.744 \\
5.4-4 & 5.1.1-2 & 0.743 \\
5.4-4 & 5.1.3-2 & 0.743 \\
5.4-4 & 5.1.4-3 & 0.743 \\
5.4-4 & 5.3.1-3 & 0.739 \\
5.4-4 & 5.1.2-1 & 0.735 \\
5.4-4 & 5.1.2-7 & 0.733 \\
5.4-4 & 5.2.1-4.2 & 0.731 \\
5.4-4 & 5.1.2-3 & 0.727 \\
5.4-4 & 5.2.1-3 & 0.723 \\
5.4-4 & 5.2.2-1 & 0.722 \\
5.4-4 & 5.2.2-6 & 0.722 \\
5.4-4 & 5.4.2-3 & 0.722 \\
5.4-4 & 5.2.2-2 & 0.721 \\
5.4-4 & 5.2.4-1.2 & 0.720 \\
5.4-4 & 5.2.4-2 & 0.720 \\
5.5-1 & 5.1.2-7 & 0.744 \\
5.5-1 & 5.2.3-1 & 0.742 \\
5.5-1 & 5.2.5-1 & 0.742 \\
5.5-1 & 5.4.1-1 & 0.735 \\
5.5-1 & 5.2.2-6 & 0.734 \\
5.5-1 & 5.2.2-4 & 0.733 \\
5.5-1 & 5.1.3-3 & 0.732 \\
5.5-1 & 5.3.1-3 & 0.730 \\
5.5-1 & 5.1.3-1 & 0.729 \\
5.5-1 & 5.3.1-2.2 & 0.727 \\
5.5-1 & 5.2.1-2 & 0.727 \\
5.5-1 & 5.2.3-2.1 & 0.723 \\
5.5-1 & 5.2.3-2 & 0.720 \\
5.5-2 & 5.2.5-1 & 0.768 \\
5.5-2 & 5.1.2-4 & 0.763 \\
5.5-2 & 5.2.3-1 & 0.762 \\
5.5-2 & 5.1.2-7 & 0.760 \\
5.5-2 & 5.2.2-4 & 0.758 \\
5.5-2 & 5.1.3-3 & 0.755 \\
5.5-2 & 5.2.5-2.1 & 0.750 \\
5.5-2 & 5.1.3-1 & 0.745 \\
5.5-2 & 5.2.3-2 & 0.745 \\
5.5-2 & 5.2.1-2 & 0.744 \\
5.5-2 & 5.4.1-1 & 0.739 \\
5.5-2 & 5.2.3-2.1 & 0.739 \\
5.5-2 & 5.4.1-2 & 0.737 \\
5.5-2 & 5.2.4-1.2 & 0.732 \\
5.5-2 & 5.2.1-1 & 0.731 \\
5.5-2 & 5.4.2-3 & 0.731 \\
5.5-2 & 5.1.4-4 & 0.728 \\
5.5-2 & 5.2.2-1 & 0.724 \\
5.5-2 & 5.3.1-3 & 0.724 \\
5.5-2 & 5.2.1-4.1 & 0.724 \\
5.5-2 & 5.2.2-6 & 0.722 \\
5.5-2 & 5.1.1-2.2 & 0.721 \\
5.5-2 & 5.2.5-3 & 0.721 \\
5.5-2 & 5.1.2-2 & 0.720 \\
5.5-2 & 5.4.1-1.1 & 0.720 \\
5.5-2 & 5.1.2-1 & 0.719 \\
5.5-5 & 5.2.1-4 & 0.756 \\
5.5-5 & 5.2.1-2 & 0.752 \\
5.5-5 & 5.2.5-1 & 0.752 \\
5.5-5 & 5.2.2-1 & 0.752 \\
5.5-5 & 5.1.3-1.1 & 0.752 \\
5.5-5 & 5.1.3-1 & 0.740 \\
5.5-5 & 5.3.1-2.2 & 0.738 \\
5.5-5 & 5.2.2-4 & 0.736 \\
5.5-5 & 5.4.1-1 & 0.735 \\
5.5-5 & 5.1.3-3 & 0.735 \\
5.5-5 & 5.2.1-1 & 0.733 \\
5.5-5 & 5.2.3-1 & 0.730 \\
5.5-5 & 5.3.1-3 & 0.730 \\
5.5-5 & 5.1.1-1 & 0.729 \\
5.5-5 & 5.2.1-4.2 & 0.728 \\
5.5-5 & 5.2.1-4.1 & 0.726 \\
5.5-5 & 5.1.1-2 & 0.725 \\
5.5-5 & 5.1.2-2 & 0.725 \\
5.5-5 & 5.1.2-1 & 0.725 \\
5.5-5 & 5.2.3-2 & 0.725 \\
5.5-5 & 5.2.2-2 & 0.721 \\
5.5-5 & 5.1.2-6 & 0.719 \\
5.5-6 & 5.2.1-4 & 0.743 \\
5.5-6 & 5.1.3-3 & 0.736 \\
5.5-6 & 5.1.4-4 & 0.728 \\
5.5-6 & 5.2.2-6 & 0.725 \\
5.5-6 & 5.2.1-2 & 0.725 \\
5.5-6 & 5.2.1-4.2 & 0.723 \\
5.5-6 & 5.5.1-1 & 0.721 \\
5.5-6 & 5.2.5-1 & 0.720 \\
5.5-7 & 5.2.1-4 & 0.754 \\
5.5-7 & 5.2.1-4.2 & 0.752 \\
5.5-7 & 5.1.3-1 & 0.728 \\
5.5-7 & 5.2.2-4 & 0.725 \\
5.5-7 & 5.2.5-1 & 0.722 \\
5.5-7 & 5.1.3-3 & 0.719 \\
5.5-8 & 5.2.3-1 & 0.789 \\
5.5-8 & 5.1.3-1 & 0.777 \\
5.5-8 & 5.2.1-2 & 0.775 \\
5.5-8 & 5.1.3-1.1 & 0.773 \\
5.5-8 & 5.2.2-4 & 0.763 \\
5.5-8 & 5.4.1-1 & 0.762 \\
5.5-8 & 5.2.1-1 & 0.759 \\
5.5-8 & 5.2.5-1 & 0.757 \\
5.5-8 & 5.2.1-4 & 0.754 \\
5.5-8 & 5.4.1-1.1 & 0.745 \\
5.5-8 & 5.1.3-3 & 0.744 \\
5.5-8 & 5.3.1-3 & 0.741 \\
5.5-8 & 5.1.2-7 & 0.741 \\
5.5-8 & 5.5.1-1 & 0.737 \\
5.5-8 & 5.2.2-1 & 0.736 \\
5.5-8 & 5.1.3-2 & 0.735 \\
5.5-8 & 5.1.2-1 & 0.734 \\
5.5-8 & 5.1.1-2.2 & 0.734 \\
5.5-8 & 5.1.1-1 & 0.734 \\
5.5-8 & 5.3.1-2.2 & 0.733 \\
5.5-8 & 5.1.1-2 & 0.733 \\
5.5-8 & 5.2.1-4.2 & 0.732 \\
5.5-8 & 5.1.2-3 & 0.732 \\
5.5-8 & 5.4.2-3 & 0.732 \\
5.5-8 & 5.1.2-4 & 0.731 \\
5.5-8 & 5.2.2-5 & 0.729 \\
5.5-8 & 5.2.3-2 & 0.728 \\
5.5-8 & 5.2.3-2.1 & 0.728 \\
5.5-8 & 5.2.2-6 & 0.727 \\
5.5-8 & 5.2.4-1.1 & 0.727 \\
5.5-8 & 5.2.1-4.1 & 0.725 \\
5.5-8 & 5.1.2-2 & 0.725 \\
5.5-8 & 5.2.5-2.1 & 0.725 \\
5.5-8 & 5.4.1-2 & 0.722 \\
5.5-8 & 5.4.2-4 & 0.721 \\
5.5-8 & 5.4.2-1 & 0.721 \\
5.5-8 & 5.2.2-3 & 0.719 \\
5.6-1 & 5.1.3-1.2 & 0.747 \\
5.6-1 & 5.2.2-4 & 0.735 \\
5.6-1 & 5.2.5-1 & 0.726 \\
5.6-2 & 5.2.1-4 & 0.754 \\
5.6-2 & 5.2.2-4 & 0.750 \\
5.6-2 & 5.2.1-1 & 0.734 \\
5.6-2 & 5.2.5-1 & 0.734 \\
5.6-2 & 5.1.3-1.2 & 0.732 \\
5.6-2 & 5.1.3-1 & 0.732 \\
5.6-2 & 5.2.1-4.2 & 0.730 \\
5.6-2 & 5.3.1-2.2 & 0.730 \\
5.6-2 & 5.2.5-4 & 0.729 \\
5.6-2 & 5.1.2-7 & 0.724 \\
5.6-2 & 5.1.3-2 & 0.723 \\
5.6-2 & 5.3.1-1 & 0.723 \\
5.6-2 & 5.2.2-1 & 0.723 \\
5.6-2 & 5.2.3-1 & 0.722 \\
5.6-2 & 5.4.2-3 & 0.721 \\
5.6-2 & 5.2.1-4.1 & 0.721 \\
5.6-2 & 5.1.2-1 & 0.720 \\
5.6-3 & 5.2.2-4 & 0.729 \\
5.6-3 & 5.1.3-1.2 & 0.725 \\
5.6-3 & 5.2.5-1 & 0.725 \\
5.6-4A & 5.2.5-1 & 0.760 \\
5.6-4A & 5.2.1-4 & 0.754 \\
5.6-4A & 5.2.2-4 & 0.748 \\
5.6-4A & 5.2.1-2 & 0.743 \\
5.6-4A & 5.2.3-1 & 0.741 \\
5.6-4A & 5.2.2-1 & 0.740 \\
5.6-4A & 5.1.3-1 & 0.736 \\
5.6-4A & 5.4.1-1 & 0.734 \\
5.6-4A & 5.1.2-7 & 0.729 \\
5.6-4A & 5.1.2-2 & 0.727 \\
5.6-4A & 5.2.3-2 & 0.727 \\
5.6-4A & 5.2.2-3 & 0.726 \\
5.6-4A & 5.2.1-4.1 & 0.726 \\
5.6-4A & 5.1.2-4 & 0.726 \\
5.6-4A & 5.4.1-1.1 & 0.724 \\
5.6-4A & 5.2.2-2 & 0.724 \\
5.6-4A & 5.1.2-3 & 0.723 \\
5.6-4A & 5.1.1-2.2 & 0.722 \\
5.6-4A & 5.4.2-3 & 0.722 \\
5.6-4A & 5.1.3-1.1 & 0.720 \\
5.6-4A & 5.2.3-2.1 & 0.719 \\
5.6-5 & 5.1.3-1.2 & 0.755 \\
5.6-5 & 5.2.3-1 & 0.729 \\
5.6-5 & 5.2.2-4 & 0.722 \\
5.6-5 & 5.2.1-1 & 0.720 \\
5.6-7 & 5.1.2-6 & 0.768 \\
5.6-7 & 5.1.3-1.2 & 0.737 \\
5.6-7 & 5.2.2-1 & 0.728 \\
5.6-7 & 5.2.2-3 & 0.725 \\
5.6-7 & 5.2.1-2 & 0.722 \\
5.6-8 & 5.2.2-1 & 0.769 \\
5.6-8 & 5.2.1-2 & 0.746 \\
5.6-8 & 5.2.1-4 & 0.743 \\
5.6-8 & 5.2.2-3 & 0.733 \\
5.6-8 & 5.1.3-1 & 0.729 \\
5.6-8 & 5.2.3-1 & 0.726 \\
5.6-8 & 5.2.5-1 & 0.724 \\
5.6-8 & 5.1.2-2 & 0.722 \\
5.6-8 & 5.2.2-2 & 0.722 \\
5.6-8 & 5.1.1-2.2 & 0.721 \\
5.6-8 & 5.2.2-4 & 0.720 \\
5.6-9 & 5.2.3-1 & 0.821 \\
5.6-9 & 5.2.1-2 & 0.802 \\
5.6-9 & 5.2.5-1 & 0.784 \\
5.6-9 & 5.1.3-1 & 0.776 \\
5.6-9 & 5.2.2-4 & 0.772 \\
5.6-9 & 5.1.1-2.2 & 0.770 \\
5.6-9 & 5.4.1-1 & 0.770 \\
5.6-9 & 5.1.2-7 & 0.769 \\
5.6-9 & 5.4.1-2 & 0.767 \\
5.6-9 & 5.2.4-1 & 0.760 \\
5.6-9 & 5.1.2-4 & 0.753 \\
5.6-9 & 5.2.1-4.1 & 0.752 \\
5.6-9 & 5.1.3-1.1 & 0.752 \\
5.6-9 & 5.3.1-3 & 0.751 \\
5.6-9 & 5.1.2-3 & 0.750 \\
5.6-9 & 5.4.1-1.1 & 0.750 \\
5.6-9 & 5.2.1-1 & 0.749 \\
5.6-9 & 5.2.3-2 & 0.744 \\
5.6-9 & 5.2.3-2.1 & 0.743 \\
5.6-9 & 5.2.2-6 & 0.741 \\
5.6-9 & 5.2.4-3 & 0.741 \\
5.6-9 & 5.4.2-3 & 0.740 \\
5.6-9 & 5.2.4-1.1 & 0.738 \\
5.6-9 & 5.2.1-4.2 & 0.734 \\
5.6-9 & 5.2.2-3 & 0.733 \\
5.6-9 & 5.2.2-1 & 0.731 \\
5.6-9 & 5.1.2-2 & 0.730 \\
5.6-9 & 5.2.5-3 & 0.729 \\
5.6-9 & 5.1.4-2 & 0.728 \\
5.6-9 & 5.1.2-1 & 0.727 \\
5.6-9 & 5.2.2-2 & 0.726 \\
5.6-9 & 5.2.4-2 & 0.726 \\
5.6-9 & 5.1.1-2 & 0.725 \\
5.6-9 & 5.1.3-4 & 0.725 \\
5.6-9 & 5.2.5-4.1 & 0.722 \\
5.6-9 & 5.1.4-4 & 0.722 \\
5.6-9 & 5.4.2-4 & 0.721 \\
5.6-9 & 5.2.5-2.1 & 0.720 \\
5.6-9 & 5.2.4-1.2 & 0.719 \\
5.6-9 & 5.1.3-1.2 & 0.719 \\
5.7-1 & 5.2.3-1 & 0.746 \\
5.7-1 & 5.1.3-3 & 0.730 \\
5.7-1 & 5.1.2-2 & 0.728 \\
5.7-1 & 5.2.1-2 & 0.724 \\
5.7-2 & 5.3.1-3 & 0.757 \\
5.7-2 & 5.4.1-1 & 0.751 \\
5.7-2 & 5.4.1-2 & 0.745 \\
5.7-2 & 5.2.1-3.1 & 0.741 \\
5.7-2 & 5.4.2-3 & 0.739 \\
5.7-2 & 5.1.3-2 & 0.738 \\
5.7-2 & 5.3.1-2.2 & 0.738 \\
5.7-2 & 5.2.1-4 & 0.737 \\
5.7-2 & 5.4.1-1.1 & 0.734 \\
5.7-2 & 5.1.2-2 & 0.731 \\
5.7-2 & 5.2.1-4.1 & 0.731 \\
5.7-2 & 5.1.3-1 & 0.729 \\
5.7-2 & 5.4.2-4 & 0.728 \\
5.7-2 & 5.1.3-3 & 0.725 \\
5.7-2 & 5.2.3-2.1 & 0.725 \\
5.7-2 & 5.2.3-1 & 0.723 \\
5.7-2 & 5.2.3-2 & 0.723 \\
5.7-2 & 5.1.3-4 & 0.722 \\
5.7-2 & 5.2.1-3 & 0.721 \\
5.7-2 & 5.1.2-4 & 0.721 \\
5.7-2 & 5.2.1-2 & 0.720 \\
5.8-1 & 5.2.1-4.2 & 0.778 \\
5.8-1 & 5.2.1-4 & 0.769 \\
5.8-1 & 5.2.2-6 & 0.755 \\
5.8-1 & 5.3.1-1 & 0.742 \\
5.8-1 & 5.2.2-4 & 0.739 \\
5.8-1 & 5.1.3-3 & 0.735 \\
5.8-1 & 5.2.1-2 & 0.734 \\
5.8-1 & 5.2.5-1 & 0.734 \\
5.8-1 & 5.1.3-1 & 0.732 \\
5.8-1 & 5.2.3-1 & 0.732 \\
5.8-1 & 5.1.2-7 & 0.729 \\
5.8-1 & 5.3.1-3 & 0.728 \\
5.8-1 & 5.4.1-1 & 0.728 \\
5.8-1 & 5.2.1-1 & 0.728 \\
5.8-1 & 5.5.1-1 & 0.725 \\
5.8-1 & 5.2.2-1 & 0.722 \\
5.8-1 & 5.2.1-4.1 & 0.720 \\
5.8-2 & 5.2.1-4.2 & 0.782 \\
5.8-2 & 5.2.1-4 & 0.782 \\
5.8-2 & 5.2.2-6 & 0.756 \\
5.8-2 & 5.3.1-1 & 0.745 \\
5.8-2 & 5.2.2-4 & 0.745 \\
5.8-2 & 5.2.5-1 & 0.745 \\
5.8-2 & 5.2.3-1 & 0.741 \\
5.8-2 & 5.2.1-1 & 0.740 \\
5.8-2 & 5.2.1-2 & 0.740 \\
5.8-2 & 5.1.3-1 & 0.739 \\
5.8-2 & 5.4.1-1 & 0.739 \\
5.8-2 & 5.1.3-3 & 0.736 \\
5.8-2 & 5.1.2-7 & 0.734 \\
5.8-2 & 5.3.1-3 & 0.733 \\
5.8-2 & 5.3.1-2.2 & 0.728 \\
5.8-2 & 5.1.2-5 & 0.728 \\
5.8-2 & 5.2.2-1 & 0.727 \\
5.8-2 & 5.2.1-4.1 & 0.726 \\
5.8-2 & 5.5.1-1 & 0.726 \\
5.8-2 & 5.1.3-1.1 & 0.721 \\
5.8-2 & 5.1.1-2 & 0.720 \\
5.8-3 & 5.3.1-1 & 0.793 \\
5.8-3 & 5.2.2-4 & 0.769 \\
5.8-3 & 5.3.1-2.2 & 0.761 \\
5.8-3 & 5.2.1-1 & 0.760 \\
5.8-3 & 5.2.4-1.1 & 0.760 \\
5.8-3 & 5.2.3-2.2 & 0.759 \\
5.8-3 & 5.3.1-2 & 0.758 \\
5.8-3 & 5.1.2-3 & 0.755 \\
5.8-3 & 5.2.4-1 & 0.749 \\
5.8-3 & 5.2.1-4 & 0.745 \\
5.8-3 & 5.1.4-2 & 0.740 \\
5.8-3 & 5.4.2-3 & 0.740 \\
5.8-3 & 5.2.4-2 & 0.740 \\
5.8-3 & 5.1.4-5 & 0.737 \\
5.8-3 & 5.2.2-6 & 0.735 \\
5.8-3 & 5.1.2-4 & 0.735 \\
5.8-3 & 5.1.2-7 & 0.734 \\
5.8-3 & 5.4.1-1 & 0.734 \\
5.8-3 & 5.1.4-1 & 0.734 \\
5.8-3 & 5.1.3-2 & 0.728 \\
5.8-3 & 5.2.3-1 & 0.727 \\
5.8-3 & 5.1.2-1 & 0.726 \\
5.8-3 & 5.2.5-1 & 0.725 \\
5.8-3 & 5.1.3-1 & 0.724 \\
5.8-3 & 5.1.3-3 & 0.724 \\
5.8-3 & 5.2.3-2.1 & 0.723 \\
5.8-3 & 5.2.3-4 & 0.722 \\
5.8-3 & 5.2.4-2.1 & 0.722 \\
5.8-3 & 5.2.3-2 & 0.721 \\
5.9-1 & 5.1.2-2 & 0.750 \\
5.9-1 & 5.2.1-3 & 0.720 \\
\end{longtable}

